\section{Limit Theorems and Methods}
\subsection{Uniqueness and Existence}
\subsubsection{For real sequence}
\bit
\item A limit inferior must exist uniquely in $[-\infty,\infty)$.
\item A limit superior must exist uniquely in $(-\infty,\infty]$.
\item If a limit exists, it is unique, and the limit inferior and limit superior are it.
\eit
\subsubsection{For real series}
If a limit exists, it is unique.
\subsubsection{For real-valued net}
\bit
\item A limit inferior must exist uniquely in $[-\infty,\infty)$, and all cluster points are greater than or equal to it.
\item A limit superior must exist uniquely in $(-\infty,\infty]$, and all cluster points are less than or equal to it.
\item If a limit exists, it is unique, and the limit inferior, limit superior, and only cluster point are it.
\eit
\subsubsection{For real-domain function at a point}
\bit
\item If an one-sided limit exists, it is unique.
\item If a limit exists, it is unique, and the left-hand limit and right-hand limit are it.
\eit
\subsubsection{For real-valued function from topological space}
\bit
\item A limit inferior must exist uniquely in $[-\infty,\infty)$.
\item A limit superior must exist uniquely in $(-\infty,\infty]$.
\item If a limit exists, it is unique, and the limit inferior and limit superior are it.
\eit
\subsubsection{For net of functions}
\bit
\item If a pointwise limit function exists, it is unique.
\item If a uniform limit function exists, it is unique, and the pointwise limit function is it.
\eit
\subsection{Limit Laws}
If the limit (including one-sided ones for functions with real domains) of sequences, series, nets, or functions with real domains, hereinafter referred to as a function, in one side exists in $\mathbb{R}\cup\{-\infty,\infty\}$, then the limit in the other side exists in $\mathbb{R}\cup\{-\infty,\infty\}$ and follows the following law.
\subsubsection{Linearity, or sum law, difference law, and constant multiple law}
The limit of a sum of constant multiples of functions is the sum of the constant multiples of the limits of the functions.
\subsubsection{Product law}
The limit of a product of functions is the product of the limits of the functions.
\subsubsection{Quotient law}
The limit of a quotient of functions is the quotient of the limits of the functions, provided that the limit of the denominator is not 0, in which the reciprocals of $\infty$ and $-\infty$ are $0$.

Thus, if $\lim f=\pm\infty$ and $\lim(fg)\in\bbR$, then $\lim g=0$.
\subsubsection{Power law}
The limit of the $n$th power of a function, in which $n$ is a positive integer, is the $n$th power of the limit of the function.
\subsubsection{Root law}
The limit of the $n$th root of a function, in which $n$ is a positive integer, is the $n$th root of the limit of the function.
\subsection{Limit inferior and limit superior laws}
\subsubsection{Addition law}
The limit inferior of a sum of functions is greater than or equal to the sum of the limit inferiors of the functions.

The limit superior of a sum of functions is less than or equal to the sum of the limit superiors of the functions.
\subsubsection{Constant multiple law}
The limit inferior of a constant multiple of a function is the constant multiple of the limit inferior of the function.

The limit superior of a constant multiple of a function is the constant multiple of the limit superior of the function.
\subsubsection{Multiplication law}
The limit inferior of a product of nonnegative real-valued functions is the product of the limit inferiors of the functions.

The limit superior of a product of nonnegative real-valued functions is the product of the limit superiors of the functions.
\subsection{Preservation of equal to}
\subsubsection{For real sequence}
Given real sequences $\langle a_n\rangle$ and $\langle b_n\rangle$ which for all $k$ that is greater than or equal to a specific $j\in \mathbb{N}$, $a_k=b_k$. If $\lim_{n\to\infty}a_n=L\in\mathbb{R}\cup\{-\infty,\infty\}$, then
\[\lim_{n\to\infty}b_n=L.\]
\subsubsection{For real series}
Given real sequences $\langle a_n\rangle$ and $\langle b_n\rangle$ which for all $k$ that is greater than or equal to a specific $j\in \mathbb{N}$:
\[\sum_{i=1}^ka_i=\sum_{i=1}^kb_i.\]
If $\sum_{i=1}^{\infty}a_i=L\in\mathbb{R}\cup\{-\infty,\infty\}$, then
\[\sum_{i=1}^{\infty}b_i=L.\]
\subsubsection{For real-valued net}
Let $A$ be a directed set. Given real-valued nets $\langle x_a\rangle_{a\in A}$ and $\langle y_a\rangle_{a\in A}$ which for all $k\in A$ such that there exists $j\in A$ with $j\leq k$, $x_k=y_k$. If $\lim_{a}x_a$ exists, then:
\[\lim_{a}x_a=\lim_{a}y_a.\]
\subsubsection{For real function}
Let $I$ be an open interval and $a\in I$, and $f(x)$ and $g(x)$ be functions defined on $I\setminus\{a\}$ which for all $x\in I\land x\neq a$, $f(x)=g(x)$. If $\lim_{x\to a}f(x)=L\in\mathbb{R}\cup\{-\infty,\infty\}$, then
\[\lim_{x\to a}g(x)=L.\]

Let $f(x)$ and $g(x)$ be functions defined on $(a,b)$ with $a<b$ which for all $x\in I$, $f(x)=g(x)$. If $\lim_{x\to a^+}f(x)=L\in\mathbb{R}\cup\{-\infty,\infty\}$, then
\[\lim_{x\to a^+}g(x)=L.\]

Let $f(x)$ and $g(x)$ be functions defined on $(b,a)$ with $a>b$ which for all $x\in I$, $f(x)=g(x)$. If $\lim_{x\to a^-}f(x)=L\in\mathbb{R}\cup\{-\infty,\infty\}$, then
\[\lim_{x\to a^-}g(x)=L.\]

Let $f(x)$ and $g(x)$ be functions defined on $(a,\infty)$ which for all $x\in I$, $f(x)=g(x)$. If $\lim_{x\to\infty}f(x)=L\in\mathbb{R}\cup\{-\infty,\infty\}$, then
\[\lim_{x\to\infty}g(x)=L.\]

Let $f(x)$ and $g(x)$ be functions defined on $(-\infty,a)$ which for all $x\in I$, $f(x)=g(x)$. If $\lim_{x\to-\infty}f(x)=L\in\mathbb{R}\cup\{-\infty,\infty\}$, then
\[\lim_{x\to-\infty}g(x)=L.\]
\subsection{Preservation of less than or equal to}
\subsubsection{For real sequence}
Given real sequences $\langle a_n\rangle$ and $\langle b_n\rangle$ which for all $k$ that is greater than or equal to a specific $j\in \mathbb{N}$, $a_k\leq b_k$. If both $\lim_{n\to\infty}a_n$ and $\lim_{n\to\infty}b_n$ are in $\mathbb{R}\cup\{-\infty,\infty\}$, then
\[\lim_{n\to\infty}a_n\leq\lim_{n\to\infty}b_n.\]
If $\lim_{n\to\infty}a_n=\infty$, then $\lim_{n\to\infty}b_n=\infty$; if $\lim_{n\to\infty}b_n=-\infty$, then $\lim_{n\to\infty}a_n=-\infty$.
\subsubsection{For real series}
Given real sequences $\langle a_n\rangle$ and $\langle b_n\rangle$ which for all $k$ that is greater than or equal to a specific $j\in \mathbb{N}$:
\[\sum_{i=1}^ka_i\leq\sum_{i=1}^kb_i.\]
If both $\sum_{i=1}^{\infty}a_i$ and $\sum_{i=1}^{\infty}b_i$ are in $\mathbb{R}\cup\{-\infty,\infty\}$, then
\[\sum_{i=1}^{\infty}a_i\leq\sum_{i=1}^{\infty}b_i.\]
If $\sum_{i=1}^{\infty}a_i=\infty$, then $\sum_{i=1}^{\infty}b_i=\infty$; if $\sum_{i=1}^{\infty}b_i=-\infty$, then $\sum_{i=1}^{\infty}a_i=-\infty$.
\subsubsection{For real-valued net}
Let $A$ be a directed set. Given real-valued nets $\langle x_a\rangle_{a\in A}$ and $\langle y_a\rangle_{a\in A}$ which for all $k\in A$ such that there exists $j\in A$ with $j\leq k$, $x_k\leq y_k$. If both $\lim_{a}x_a$ and $\lim_{a}y_a$ exist, then:
\[\lim_{a}x_a\leq\lim_{a}y_a.\]
\subsubsection{For real function}
Let $I$ be an open interval and $a\in I$, and $f(x)$ and $g(x)$ be functions defined on $I\setminus\{a\}$ which for all $x\in I\land x\neq a$, $f(x)\leq g(x)$. If both $\lim_{x\to a}f(x)$ and $\lim_{x\to a}g(x)$ are in $\mathbb{R}\setminus\{-\infty,\infty\}$, then
\[\lim_{x\to a}f(x)\leq\lim_{x\to a}g(x).\]
If $\lim_{x\to a}f(x)=\infty$, then $\lim_{x\to a}g(x)=\infty$; if $\lim_{x\to a}g(x)=-\infty$, then $\lim_{x\to a}f(x)=-\infty$.

Let $f(x)$ and $g(x)$ be functions defined on $(a,b)$ with $a<b$ which for all $x\in I$, $f(x)\leq g(x)$. If both $\lim_{x\to a^+}f(x)$ and $\lim_{x\to a^+}g(x)$ are in $\mathbb{R}\cup\{-\infty,\infty\}$, then
\[\lim_{x\to a^+}f(x)\leq\lim_{x\to a^+}g(x).\]
If $\lim_{x\to a^+}f(x)=\infty$, then $\lim_{x\to a^+}g(x)=\infty$; if $\lim_{x\to a^+}g(x)=-\infty$, then $\lim_{x\to a^+}f(x)=-\infty$.

Let $f(x)$ and $g(x)$ be functions defined on $(b,a)$ with $a>b$ which for all $x\in I$, $f(x)\leq g(x)$. If both $\lim_{x\to a^-}f(x)$ and $\lim_{x\to a^-}g(x)$ are in $\mathbb{R}\cup\{-\infty,\infty\}$, then
\[\lim_{x\to a^-}f(x)\leq\lim_{x\to a^-}g(x).\]
If $\lim_{x\to a^-}f(x)=\infty$, then $\lim_{x\to a^-}g(x)=\infty$; if $\lim_{x\to a^-}g(x)=-\infty$, then $\lim_{x\to a^-}f(x)=-\infty$.

Let $f(x)$ and $g(x)$ be functions defined on $(a,\infty)$ which for all $x\in I$, $f(x)\leq g(x)$. If both $\lim_{x\to\infty}f(x)$ and $\lim_{x\to\infty}g(x)$ are in $\mathbb{R}\cup\{-\infty,\infty\}$, then
\[\lim_{x\to\infty}f(x)\leq\lim_{x\to\infty}g(x).\]
If $\lim_{x\to\infty}f(x)=\infty$, then $\lim_{x\to\infty}g(x)=\infty$; if $\lim_{x\to\infty}g(x)=-\infty$, then $\lim_{x\to\infty}f(x)=-\infty$.

Let $f(x)$ and $g(x)$ be functions defined on $(-\infty,a)$ which for all $x\in I$, $f(x)\leq g(x)$. If both $\lim_{x\to-\infty}f(x)$ and $\lim_{x\to-\infty}g(x)$ are in $\mathbb{R}\cup\{-\infty,\infty\}$, then
\[\lim_{x\to-\infty}f(x)\leq\lim_{x\to-\infty}g(x).\]
If $\lim_{x\to-\infty}f(x)=\infty$, then $\lim_{x\to-\infty}g(x)=\infty$; if $\lim_{x\to-\infty}g(x)=-\infty$, then $\lim_{x\to-\infty}f(x)=-\infty$.
\subsection{Squeeze (夾擠) theorem or Sandwich (三明治) theorem}
\subsubsection{For real sequence}
Given real sequences $\langle a_n\rangle$, $\langle b_n\rangle$, and $\langle c_n\rangle$ which for all $k$ that is greater than or equal to a specific $j\in \mathbb{N}$:
\[a_k\leq c_k\leq b_k.\]
If
\[\lim_{n\to\infty}a_n=\lim_{n\to\infty}b_n=L\in\mathbb{R}\cup\{-\infty,\infty\},\]
then
\[\lim_{n\to\infty}c_n=L.\]
\subsubsection{For real series}
Given real sequences $\langle a_n\rangle$, $\langle b_n\rangle$, and $\langle c_n\rangle$ which for all $k$ that is greater than or equal to a specific $j\in \mathbb{N}$:
\[\sum_{i=1}^ka_i\leq\sum_{i=1}^kc_i\leq\sum_{i=1}^kb_i.\]
If
\[\sum_{i=1}^{\infty}a_i=\sum_{i=1}^{\infty}b_i=L\in\mathbb{R}\cup\{-\infty,\infty\},\]
then: 
\[\sum_{i=1}^{\infty}c_i=L.\]
\subsubsection{For real-valued net}
Let $A$ be a directed set. Given real-valued nets $\langle x_a\rangle_{a\in A}$, $\langle y_a\rangle_{a\in A}$, and $\langle z_a\rangle_{a\in A}$ which for all $k\in A$ such that there exists $j\in A$ with $j\leq k$:
\[x_k\leq z_k\leq y_k.\]
If
\[\lim_{a}x_a=\lim_{a}y_a=L,\]
then: 
\[\lim_{a}z_a=L.\]
\begin{proof}
Start from the limits of $x_a$ and $y_a$, by definition of limit:

For every $\varepsilon\in\mathbb{R}_{>0}$, there exists some $a_1\in A$ such that for every $b\in A$ with $b\geq a_1$, the point $L-\varepsilon\leq x_b\leq L+\varepsilon$.

For every $\varepsilon\in\mathbb{R}_{>0}$, there exists some $a_2\in A$ such that for every $b\in A$ with $b\geq a_2$, the point $L-\varepsilon\leq y_b\leq L+\varepsilon$.

Choose $a\in A$ such that $a_1\leq a$ and $a_2\leq a$.

so,
\[L-\varepsilon\leq x_{a}\leq z_{a}\leq y_{a}\leq L+\varepsilon.\]

Since $\varepsilon > 0$ was arbitrary, by definition of limit:
\[\lim_{a}z_a=L.\]
\end{proof}
\subsubsection{For real function}
Let $I$ be an open interval and $a\in I$, and $f(x)$, $g(x)$, and $h(x)$ be functions defined on $I\setminus\{a\}$ which for all $x\in I\land x\neq a$, $f(x)\leq h(x)\leq g(x)$. If
\[\lim_{x\to a}f(x)=\lim_{x\to a}g(x)=L\in\mathbb{R}\setminus\{-\infty,\infty\},\]
then
\[\lim_{x\to a}h(x)=L.\]

Let $f(x)$, $g(x)$, and $h(x)$ be functions defined on $(a,b)$ with $a<b$ which for all $x\in I$, $f(x)\leq h(x)\leq g(x)$. If
\[\lim_{x\to a^+}f(x)=\lim_{x\to a^+}g(x)=L\in\mathbb{R}\setminus\{-\infty,\infty\},\]
then
\[\lim_{x\to a^+}h(x)=L.\]

Let $f(x)$, $g(x)$, and $h(x)$ be functions defined on $(b,a)$ with $a>b$ which for all $x\in I$, $f(x)\leq h(x)\leq g(x)$. If
\[\lim_{x\to a^-}f(x)=\lim_{x\to a^-}g(x)=L\in\mathbb{R}\setminus\{-\infty,\infty\},\]
then
\[\lim_{x\to a^-}h(x)=L.\]

Let $f(x)$, $g(x)$, and $h(x)$ be functions defined on $(a,\infty)$ which for all $x\in I$, $f(x)\leq h(x)\leq g(x)$. If
\[\lim_{x\to\infty}f(x)=\lim_{x\to\infty}g(x)=L\in\mathbb{R}\setminus\{-\infty,\infty\},\]
then
\[\lim_{x\to\infty}h(x)=L.\]

Let $f(x)$, $g(x)$, and $h(x)$ be functions defined on $(-\infty,a)$ which for all $x\in I$, $f(x)\leq h(x)\leq g(x)$. If
\[\lim_{x\to-\infty}f(x)=\lim_{x\to-\infty}g(x)=L\in\mathbb{R}\setminus\{-\infty,\infty\},\]
then
\[\lim_{x\to-\infty}h(x)=L.\]
\begin{proof}
Take the first statement for example. The other statements can be proven similarly.

Start from the limits of $f$ and $g$, by definition of limit:
\[\forall\varepsilon>0, \exists\delta_1>0\text{\ s.t.\ }0<|x-a|<\delta_1\implies|f(x)-L|<\varepsilon,\]
and
\[\forall\varepsilon>0, \exists\delta_2>0\text{\ s.t.\ }0<|x-a|<\delta_2\implies|g(x)-L|<\varepsilon,\]

Choose:
\[\delta\coloneq\min(\delta_1,\delta_2),\]

so,
\[L-\varepsilon<f(x)\leq h(x)\leq g(x)<L+\varepsilon.\]

Since $\varepsilon > 0$ was arbitrary, by definition of limit:
\[\lim_{x\to a}h(x)=L.\]
\end{proof}
\subsection{Limit of real-valued net theorems}
\subsubsection{Continuity principle}
If real-valued net $\langle x_a\rangle_{a\in A}$ with $\lim_ax_a=L\in\bbR$, then for any real-valued function $f$ continuous at $L$,
\[\lim_af(x_a)=f(L).\]
If real-valued net $\langle x_a\rangle_{a\in A}$ with $\lim_ax_a=\pm\infty$, then for any real-valued function $f$ such that $\lim_{x\to\pm\infty}f(x)=L\in\bbR\cup\{-\infty,\infty\}$,
\[\lim_af(x_a)=L.\]
\subsubsection{Monotone Convergence Theorem (MCT) (單調收斂定理)}
For any nonempty, non-decreasing net of real numbers $\langle a_n\rangle_{n\in A}$, it is bounded-above iff it converges to a real number, and its limit is its supremum.
\begin{proof}
$\Rightarrow$ direction:

Let $\{a_n\}$ be the set of values of $\langle a_n\rangle_{n\in A}$. Assume that $\{a_n\}$ is nonempty and bounded-above by $c=\sup_na_n$.

We have
\[\forall\varepsilon>0:\,\exists M\in A\text{\ s.t.\ }c\geq a_M>c-\varepsilon,\]
since otherwise $c-\varepsilon$ is a strictly smaller upper bound of $\langle a_n\rangle$, contradicting the definition of the supremum. 

Then since $\langle a_n\rangle$ is non-decreasing, and $c$ is an upper bound:
\[\forall\varepsilon>0:\,\exists M\in A\text{\ s.t.\ }\forall n\geq M:\,|c-a_n|=c-a_n\leq c-a_M=|c-a_M|<\varepsilon.\]
$\Leftarrow$ direction:

Assume that $\lim_na_n=L\in\bbR$. Assume for contradiction that for some $M\in A$,
\[a_M>L.\]
Let
\[\varepsilon=\frac{a_M-L}{2}>0.\]
Since
\[\forall n\geq M:\, a_n\geq a_M>L,\]
we have
\[\forall n\geq M:\,|a_n-L|=(a_n-a_M)+(a_M-L)\geq a_M-L>\varepsilon.\]
This contradicts that
\[\lim_na_n=L.\]
\end{proof}
For any nonempty, non-increasing net of real numbers $\langle a_n\rangle_{n\in A}$, it is bounded-below iff it converges to a real number, and its limit is its infimum.
\subsection{Limit of real sequence theorems}
\subsubsection{Bolzano–Weierstrass Theorem (波爾查諾-魏爾斯特拉斯定理)}
Every bounded sequence in an Euclidean space $\bbR^n$ has a convergent subsequence.

\begin{proof}
    For $\bbR$, take any bounded sequence $(x_n)$ in $\bbR$. Then there exist real numbers $m,M$ such that
    \[m\leq x_n\leq M,\quad\forall n.\]
    So the sequence lies in the closed interval $[m,M]$.

    For $\varepsilon=1$. Divide $[m,M]$ into disjoint subintervals with length $\frac{M-m}{2^\varepsilon}$:
    \[\qty[m,\frac{m+M}{2}],\qty[\frac{m+M}{2},M].\]
    Since $(x_n)$ is infinite, by the pigeonhole principle, one of the subintervals must contain infinitely many terms of $(x_n)$. Denote it by $I_1$ and pick the subsequence $\qty(x_n^{(1)})$ of $(x_n)$ that is entirely in $I_1$.

    Repeat it for $\varepsilon=2,3,\ldots,$ by dividing $I_{\varepsilon-1}$ into two halves, we can construct a nested sequence of subsequences
    \[\qty(x_n^{(1)})\supseteq \qty(x_n^{(2)})\supseteq\ldots\]
    where $\qty(x_n^{(k)})$ lies in a closed interval of length $\frac{M-m}{2^k}$, $I_k$.

    Pick the diagonal sequence $y_k=x_k^{(k)}$. Then
    \[y_k\in\qty(x_n^{(k)})\subseteq I_k,\]
    and the sequence $(y_k)$ is a convergent subsequence of $(x_n)$ because
    \[\forall N\in\bbN\colon m,n\geq N\implies |y_m-y_n|\leq\frac{M-m}{2^N}.\]

    To generalized to $\bbR^n$, consider each coordinate sequence in $\bbR$, it has a convergent subsequence. Diagonalization gives a subsequence converging in all coordinates.
\end{proof}
\subsubsection{Abel transformation or summation by parts}
Let $\langle f_k\rangle$ and $\langle g_k\rangle$ be two sequences.
\[\sum_{k=m}^nf_k\qty(g_{k+1}-g_k)=\qty(f_{n+1}g_{n+1}-f_mg_m)-\sum_{k=m}^ng_{k+1}\qty(f_{k+1}-f_k).\]
\[\sum_{k=m}^nf_k\qty(g_k-g_{k-1})=\qty(f_ng_n-f_{m-1}g_{m-1})-\sum_{k=m}^ng_{k-1}\qty(f_k-f_{k-1}).\]
\[\sum_{k=m}^n\qty(f_k\qty(g_{k+1}-g_k)+g_k\qty(f_k-f_{k-1}))=f_ng_{n+1}-f_{m-1}g_m.\]
\subsection{Limit of real function theorems}
\subsubsection{Limits involving quotient functions}
Let \( a \) and \( b \) be real numbers, set
\[A=\left\{f\colon U\subseteq\mathbb{R}\to\mathbb{R} \middle | f(x) = x \lor \ln(f(x)) \in A \lor e^{f\left(x\right)}  \in A \right\},\]
and function $f\in A$. Then:
\[\lim_{x \to \infty} \frac{\left(f\left(x\right)\right)^a}{\left(f\left(x\right)\right)^b} = \infty, \quad a > b \]
\[ \lim_{x \to \infty} \frac{af\left(x\right)}{bf\left(x\right)} = \frac{a}{b}, \quad b \neq 0 \]
\[ \lim_{x \to \infty}\frac{n^{af\left(x\right)}}{bf\left(x\right)} = \infty, \quad a,b > 0 \land  n > 1 \]
\[ \lim_{x \to \infty}\frac{n^{af\left(x\right)}}{bf\left(x\right)} = 0, \quad a,b > 0 \land  0\leq n<1 \]
\[ \lim_{x \to \infty}\frac{af\left(x\right)}{b\log_n f\left(x\right)} = \infty, \quad a,b > 0 \land  n > 1 \]
\subsubsection{BV function has limit}
If function $f\colon[a,b)\to\bbR$ with $b\in(a,\infty]$ is of bounded variation on $[a,b)$, then $\lim_{x\to b^-}f(x)$ exists in $\bbR$.
\begin{proof}
PLACEHOLDER
\end{proof}